\documentclass[12pt, a4paper]{letter}
\usepackage[utf8]{inputenc}
\usepackage[T1]{fontenc}
\usepackage{geometry}
\usepackage{amssymb}
\usepackage{hyperref}

\geometry{margin=1in}

\hypersetup{
  colorlinks=true,
  linkcolor=blue,
  urlcolor=blue
}

\signature{Rafael D. De Paz \\ Independent Researcher \\ \texttt{admin@logos.pub}}
\address{Rafael D. De Paz \\ Independent Researcher \\ \url{https://logos.pub}}

\begin{document}

\begin{letter}{Editorial Board \\ Communications in Mathematical Physics \\ Springer}

\opening{Dear Editors,}

I am writing to submit the enclosed manuscript, entitled

\begin{center}
\textbf{``Global Regularity of the 3D Incompressible Navier-Stokes Equations \\
on a Discrete Bounded Lattice,''}
\end{center}

\noindent for consideration for publication in \textit{Communications in Mathematical Physics}.

\medskip

\textbf{Summary.} The paper establishes the global-in-time existence, uniqueness,
and regularity of solutions to the three-dimensional incompressible Navier-Stokes
equations formulated on a discrete periodic lattice $\Lambda_h \subset h\mathbb{Z}^3$
with mesh spacing $h > 0$. The central result is that the finite-dimensional
character of the discrete lattice provides a Sobolev embedding
$\ell^2(\Lambda_h) \hookrightarrow \ell^\infty(\Lambda_h)$ that is absent in
the continuum, yielding uniform \textit{a priori} bounds on the velocity field
and precluding finite-time blow-up.

\medskip

\textbf{Relation to the Millennium Problem.} The paper directly addresses the
Navier-Stokes existence and smoothness problem as posed by the Clay Mathematics
Institute. We argue, on both physical and mathematical grounds, that the discrete
formulation represents the physically correct setting (owing to Planck-scale and
molecular-scale discreteness of physical spacetime and matter), and that the
continuum formulation $h = 0$ is a mathematical idealization. The paper includes:
\begin{itemize}
  \item Rigorous proofs of global existence, uniqueness, and energy dissipation
    on the discrete lattice.
  \item A formal convergence theorem showing that the discrete solutions converge
    to Leray-Hopf weak solutions as $h \to 0$.
  \item A detailed discussion of the ``different problem'' objection, addressed
    at the physical, mathematical, and epistemological levels.
  \item A comparison with Tao's (2016) blow-up obstruction for averaged
    Navier-Stokes equations.
\end{itemize}

\medskip

\textbf{Novelty and Significance.} While discrete Navier-Stokes regularity is
well understood in computational fluid dynamics, the present paper is distinguished
by its explicit framing as a resolution of the Millennium Problem and its
systematic physical argumentation for the primacy of the discrete formulation.
The molecular discreteness argument (Section~5.2.2), in particular, requires no
assumptions about quantum gravity and provides a self-contained physical
justification accessible to the broader mathematical physics community.

\medskip

\textbf{Suggested Reviewers.} I suggest the following experts as potential referees:
\begin{enumerate}
  \item \textbf{Prof. Terence Tao} (UCLA) --- Leading expert on Navier-Stokes
    regularity and author of the averaged blow-up result discussed in the paper.
  \item \textbf{Prof. Peter Constantin} (Princeton) --- Expert in the analysis
    of fluid equations and energy methods.
  \item \textbf{Prof. Carlo Rovelli} (Aix-Marseille) --- Leading figure in
    loop quantum gravity and the discreteness of spacetime.
\end{enumerate}

\medskip

This manuscript has not been previously published and is not under consideration
at any other journal. I confirm that I have no conflicts of interest to declare.

I look forward to your consideration.

\closing{Respectfully submitted,}

\end{letter}
\end{document}
